\section{Laundering Topology and Heuristic Pattern Detection}
\label{sec:pattern_detection}

Modern blockchain laundering syndicates do not disburse illicit funds in random configurations; instead, their operations adhere to identifiable topological subgraphs designed to balance two competing operational objectives: \emph{capital preservation} (minimizing slippage and gas fees) and \emph{forensic obfuscation} (severing deterministic clustering linkages). In this section, we formalize the structural characteristics and algorithmic criteria for identifying four primary money laundering typologies: Peel Chains, Smurfing/Structuring, Mixer Interactions, and Cross-Chain Bridge Gateways.

\subsection{Peel Chain Topology and Change Disambiguation}

A \textbf{Peel Chain} is an elongated sequential laundering topology in which an attacker repeatedly forwards the predominant fraction of funds to a newly generated intermediate address while ``peeling off'' smaller fractional payments to external destinations (such as non-KYC exchanges, peer-to-peer OTC desks, prepaid card services, or operational cashout nodes).

\subsubsection{Formal Structural Definition}
Let a directed subpath of length $L \ge 2$ be denoted by $\mathcal{P} = (u_0, e_1, u_1, e_2, u_2, \dots, e_L, u_L)$. The path $\mathcal{P}$ constitutes an active Peel Chain if, for every intermediate hop $i \in \{1, \dots, L-1\}$, the node $u_i$ satisfies:

\begin{enumerate}
    \item \textbf{Low In-Degree and Out-Degree Parity:}
    \begin{equation}
        \text{deg}^-(u_i) = 1 \quad \text{and} \quad \text{deg}^+(u_i) = 2
    \end{equation}
    where one outgoing edge $e_i^{\text{forward}} = (u_i, u_{i+1})$ forwards the core balance, and the second edge $e_i^{\text{peel}} = (u_i, p_i)$ peels off operational expenditure to an external node $p_i \ne u_{i+1}$.
    
    \item \textbf{Asymmetric Volume Ratio:}
    \begin{equation}
        \frac{w(e_i^{\text{forward}})}{w(e_i^{\text{peel}})} \ge \Omega_{\text{peel}}, \quad \text{with } \Omega_{\text{peel}} \ge 4.0
    \end{equation}
    typically reflecting an $80\% / 20\%$ or $95\% / 5\%$ volume asymmetry.
    
    \item \textbf{Temporal Rapid Forwarding:} The dwell time $\Delta t_i = t(e_i^{\text{forward}}) - t(e_i^{\text{in}})$ between receiving incoming funds and initiating outgoing transfers satisfies:
    \begin{equation}
        \Delta t_i \le \Delta T_{\text{peel\_max}}, \quad \Delta T_{\text{peel\_max}} = 3600\,\text{s (1 hour)}
    \end{equation}
    demonstrating programmatic, automated forwarding scripts.
\end{enumerate}

\subsubsection{Change Address Disambiguation}
In account-based chains, unlike UTXO networks, there is no inherent concept of a ``change address''. However, attackers programmatically simulate this behavior by creating ephemeral EOAs. Our heuristic engine disambiguates the true continuation node $u_{i+1}$ from the peeled offramp $p_i$ via:
\begin{equation}
    u_{i+1} = \arg\max_{v \in \mathcal{N}^+(u_i)} \left\{ w(u_i, v) \cdot \mathbb{I}_{\{\text{Type}(v) = \text{EOA}\}} \cdot \exp(-\mu \cdot \text{Age}(v)) \right\}
\end{equation}
where $\text{Age}(v)$ is the transaction count of the recipient. A completely fresh address with zero prior transaction history receiving $>80\%$ of the balance is labeled with $99.2\%$ confidence as the continuation peel node.

\subsection{Smurfing and Multi-Tier Structuring}

\textbf{Smurfing} (or structuring) represents a high-entropy topological dispersion pattern wherein a high-value capital pool is split into numerous micro-transfers distributed across tens or hundreds of intermediary mule accounts, frequently to circumvent automated Anti-Money Laundering (AML) and Currency Transaction Report (CTR) monitoring thresholds.

\subsubsection{Fan-Out Dispersion Stage}
Let $u \in \mathcal{V}$ be a distributing node. Node $u$ is flagged as a \textbf{Fan-Out Smurfing Gateway} if it satisfies:
\begin{equation}
    \text{deg}^+(u) \ge K_{\text{smurf\_fan}} \quad (K_{\text{smurf\_fan}} \ge 5)
\end{equation}
and the transacted values exhibit low variance, indicating artificial structuring:
\begin{equation}
    \text{CV}(w_{\text{out}}(u)) = \frac{\sigma(w_{\text{out}}(u))}{\mu(w_{\text{out}}(u))} \le \epsilon_{\text{variance}} \quad (\epsilon_{\text{variance}} \le 0.35)
\end{equation}
where $\text{CV}$ is the coefficient of variation across outgoing edge amounts, and $\sigma, \mu$ are the standard deviation and mean, respectively.

\subsubsection{Fan-In Consolidation Stage}
Smurfed funds cannot remain permanently dispersed; to be liquidated into fiat currency, they must eventually re-aggregate at centralized exchange deposit accounts or liquidity pools. A downstream node $v_{\text{sink}}$ is classified as a \textbf{Fan-In Aggregation Hub} if:
\begin{equation}
    \text{deg}^-(v_{\text{sink}}) \ge K_{\text{smurf\_fan}} \quad \text{and} \quad \sum_{e \in \mathcal{E}^-(v_{\text{sink}})} w_e \ge \Pi_{\text{recovery\_ratio}} \cdot W_0
\end{equation}
within a bounded observation window $[t_0, t_0 + \Delta T_{\text{recon}}]$.

\subsection{Mixer Protocols and Anonymity Pools}

Smart contract-based privacy protocols represent deliberate cryptographic barriers in the transaction graph. Unlike standard peer-to-peer transfers, interactions with mixing contracts exhibit distinct topological signatures.

\subsubsection{On-Chain Mixer Identification}
A node $m \in \mathcal{V}$ is labeled as a \textbf{Mixer Contract} if its address matches our continuously updated cryptographic directory of known zero-knowledge and coin-join protocols:
\begin{equation}
    m \in \mathcal{M}_{\text{known}} = \{\text{Tornado.Cash Routers}, \text{Railgun}, \text{Cyclone}, \dots\}
\end{equation}
or if the contract bytecode matches verified zk-SNARK verification logic (such as pairings over the BN254 / alt\_bn128 curve).

\subsubsection{Fixed-Denomination Accumulator Signature}
Mixers require uniform deposit denominations to maintain an unpartitioned anonymity set:
\begin{equation}
    w(u, m) \in \{0.1, 1.0, 10.0, 100.0\}\,\text{ETH} \quad \text{or} \quad \{100, 1000, 10000, 100000\}\,\text{USDT}
\end{equation}
When an incoming wire satisfies this discrete value constraint and targets a verified mixer contract, the destination node is tagged with $\text{Risk} = 100$, and the search engine logs a severe \emph{Cryptographic Privacy Severance} event.

\subsection{Cross-Chain Bridges and Multi-Ledger Hopping}

With increased regulatory scrutiny on Ethereum mixers, sophisticated cyber groups frequently execute cross-chain hops to disrupt graph traversal algorithms constrained to a single blockchain ledger.

\subsubsection{Bridge Contract Attribution}
A destination address $b \in \mathcal{V}$ is categorized as a \textbf{Cross-Chain Bridge Gateway} if:
\begin{equation}
    b \in \mathcal{B}_{\text{bridge}} = \{\text{Stargate}, \text{Across}, \text{Hop}, \text{Wormhole}, \text{Arbitrum Bridge}, \dots\}
\end{equation}
or if the transaction input data invokes known lock-and-burn function selectors (e.g., \texttt{depositTo()}, \texttt{bridgeTokens()}, \texttt{swapAndBridge()}).

\subsubsection{Cross-Ledger State Transition Invariant}
When tainted capital reaches a bridge contract $b$ on source chain $\mathcal{C}_{\text{src}}$ at timestamp $t_{\text{lock}}$, the physical asset is locked or burned. The engine records a synthetic state transition:
\begin{equation}
    \Phi_{\text{bridge\_event}} = \langle \mathcal{C}_{\text{src}}, \mathcal{C}_{\text{dst}}, u_{\text{sender}}, b, w_{\text{locked}}, \kappa_{\text{src}}, t_{\text{lock}} \rangle
\end{equation}
This triggers an automated cross-ledger index lookup on destination network $\mathcal{C}_{\text{dst}}$ for a corresponding minting or release event within the temporal tolerance window $[t_{\text{lock}}, t_{\text{lock}} + 1800\,\text{s}]$.

\subsection{Composite Multi-Factor Heuristic Risk Scoring}

To provide investigators with immediate actionable intelligence, our framework computes a composite node risk score $\text{Risk}(u) \in [0, 100]$ synthesizing topological, behavioural, and attribution signals:

\begin{equation}
    \text{Risk}(u) = \min\left(100, \; \sum_{k=1}^{7} \omega_k \cdot \mathcal{I}_k(u) \right)
    \label{eq:risk_score}
\end{equation}
where $\mathcal{I}_k(u)$ represents heuristic indicator functions defined in Table~\ref{tab:risk_indicators}.

\begin{table}[h]
\centering
\caption{Heuristic Risk Scoring Indicator Weights.}
\label{tab:risk_indicators}
\begin{tabular}{clcc}
\toprule
\textbf{Indicator} & \textbf{Forensic Criteria} & \textbf{Weight ($\omega_k$)} & \textbf{Severity} \\
\midrule
$\mathcal{I}_1(u)$ & Direct Interaction with Verified Sanctioned/OFAC Entity & 100 & Critical \\
$\mathcal{I}_2(u)$ & Direct Deposit to Zero-Knowledge Mixer (e.g., Tornado) & 95 & Critical \\
$\mathcal{I}_3(u)$ & Participation in Active Peel Chain ($\ge 3$ sequential hops) & 75 & High \\
$\mathcal{I}_4(u)$ & High-Fanout Smurfing Gateway ($\text{deg}^+ \ge 10$) & 70 & High \\
$\mathcal{I}_5(u)$ & Cross-Chain Bridge Lock Event & 65 & Medium \\
$\mathcal{I}_6(u)$ & Interaction with Unregulated / Instant Non-KYC Swap & 60 & Medium \\
$\mathcal{I}_7(u)$ & Taint Dilution Coefficient $\tau(u) \ge 0.5$ & 50 & Medium \\
\bottomrule
\end{tabular}
\end{table}

This multi-factor formulation ensures that addresses directly involved in high-risk obfuscation vectors are immediately escalated to the top of the forensic priority queue.
